\section{Discussion and Limitations}
\label{sec:discussion}

We are deliberately explicit about what this instrument is and is not.

\paragraph{The evaluation loop is closed.}
The same model family generates items, answers them, and judges the
answers. Any systematic blind spot---a regulatory nuance the model
family consistently misreads---is invisible to the loop and can poison
both the gold standard and its grading. The real-FDA track partly
mitigates this by injecting externally authored content, but the judge is
still an LLM. External validity against human expert judgment is unproven.

\paragraph{Synthetic examinees are correlated, not independent.}
Classical IRT calibration assumes many independent test-takers
(hundreds to thousands). Our 45-cell grid is three base models under
prompt and temperature variation; cells sharing a base model are highly
correlated. This makes the difficulty estimates more informative than a
two-profile design but less trustworthy than the examinee count alone
would suggest. Standard errors should be read in that light.

\paragraph{Difficulty uncertainty is currently analytic.}
The $b \pm \mathrm{SE}$ reported in Section~\ref{sec:results} uses an
analytic single-item logit SE. The intended final uncertainty is the
posterior standard deviation from a variational Rasch fit
(\texttt{py-irt}), which is implemented but requires regenerating the raw
response matrix. We expect the posterior SDs to be \emph{wider} than the
analytic values, given only 45 examinees---which is the honest message,
not a defect to hide. \todo{Report the py-irt posterior SDs once
regenerated and compare to the analytic SE.}

\paragraph{Judge error is real and measured.}
Re-grading introduced a measurable judge-error rate (failed judge calls
mapped to a failing score). We treat this as part of the measurement
noise rather than assuming it away, but a human-validated subset would be
needed to quantify judge bias as opposed to judge variance.

\paragraph{What would make this credible as measurement.}
Three things, in priority order: (i)~\emph{human expert validation} of a
stratified sample of gold standards and verdicts, providing an external
anchor for the difficulty scale; (ii)~a formal \emph{real-vs-synthetic}
test (Mann--Whitney~$U$) on the difficulty gap reported descriptively in
Section~\ref{sec:results}; and (iii)~\emph{cross-model} examinee
expansion to reduce within-family correlation. The Red Herring mandate is
enforced at generation but not audited---no automated check verifies that
a genuine Red Herring is present and correctly explained in the gold
standard.

\paragraph{Honest framing of the contribution.}
This is a pipeline and an instrument, not a psychometric certification.
The $\Delta\theta$ decomposition (Section~\ref{sec:deltatheta}) is the
most defensible empirical result: it is a within-model comparison that is
invariant to the scale's absolute anchoring, and it delivers a concrete,
falsifiable claim---system-prompt strictness mostly buys citation format.
